% QDTS paper skeleton: section and subsection headers only.
% Sources: QDTS_paper_writing_plan.md (outline, page budgets, figure/table IDs, claims C1-C14)
%          EVALUATION_FRAMEWORK_SECTION.md (merged into Section III)
% Target: IEEE ICDE 2028, regular research track. At most 12 pages excluding references and the
% AI-content acknowledgement. No appendix: anything marked [TR] goes to the arXiv technical report.
%
% Page budget per section (ICDE column of the plan):
%   I 1.0 | II 0.5 | III 1.0 | IV 0.75 | V 2.0 | VI 1.0 | VII 1.25 | VIII 1.5 | IX 1.25 | X 1.0
%   XI 0.5 | XII 0.25 | total 12.0
% Section VI holds the evaluation fairness material and the information-access table (T2), so it
% takes 0.25 pages from Section III.
%
% Rule between Section III (Preliminaries) and Section VI (Experimental Setup): Section III holds
% definitions, no parameter values. Section VI holds the values and the fairness rules, and refers
% back to Section III instead of re-defining anything. Runtime and hardware belong in Section XI.

\documentclass[conference]{IEEEtran}

\usepackage{amsmath,amssymb,amsthm}
\usepackage{graphicx}
\usepackage{booktabs}
\usepackage{cite}
\usepackage{xcolor}
\usepackage{hyperref}

\newtheorem{lemma}{Lemma}
\newtheorem{proposition}{Proposition}
\newtheorem{theorem}{Theorem}

\newcommand{\todo}[1]{\textcolor{red}{[TODO: #1]}}
\newcommand{\planned}[1]{\textcolor{blue}{[PLANNED: #1]}}

\begin{document}

% Candidate titles (choose after student go/no-go 2, plan section E):
%   1. Query-Driven Trajectory Simplification Is Nearly Separable, and Therefore Learnable
%   2. What Matters in Query-Driven Trajectory Simplification? Where the Budget Goes, Then Which
%      Points Stay
%   3. Learning to Spend a Global Point Budget: Teacher-Guided Query-Driven Trajectory
%      Simplification
\title{\todo{Title}}

\author{\IEEEauthorblockN{\todo{Authors}}
\IEEEauthorblockA{\todo{Affiliation}}}

\maketitle

\begin{abstract}
% Seven sentences (plan section E): context, gap, approach, structure, findings with numbers,
% student, implication.
\todo{Abstract}
\end{abstract}

\begin{IEEEkeywords}
\todo{Keywords}
\end{IEEEkeywords}

% ---------------------------------------------------------------------------------------------
\section{Introduction}
\label{sec:intro}
% 1.0 page. Draft: INTRODUCTION_DRAFT.md. No subsections. Paragraph flow:
%   motivation (offline, once-per-database compression under one global budget);
%   error-driven simplification and its limits; QDTS (RL4QDTS, MLSimp), crediting RL4QDTS's
%   "Issue 1: uniform compression ratio" and MLSimp's global-vs-per-trajectory result;
%   Example 1 with Fig. 1; three challenges; teacher-student framework with marginality and slack;
%   how each challenge is answered; Example 2; findings (three-number answer with budget
%   condition); contributions; limitations paragraph; outline.
% Figure: F1 (page-1 map teaser).

% ---------------------------------------------------------------------------------------------
\section{Related Work}
\label{sec:related}
% 0.5 page.

\subsection{Error-Driven Trajectory Simplification}
\label{sec:related-edts}
% DP/RDP, Bottom-Up, Span-Search, SED/PED/DAD/SAD, RLTS, S3.

\subsection{Query-Driven Trajectory Simplification}
\label{sec:related-qdts}
% RL4QDTS and MLSimp; Zhang et al. as the evaluation precedent. What is new here: the quantified,
% order-free split, the separability result, and that naive global allocation is worse than equal
% budgets.

\subsection{Attribution and Workload-Aware Optimisation}
\label{sec:related-attribution}
% Ablation, factorial and Shapley attribution; workload-aware DB optimisers (Flood, Qd-tree).

\subsection{Learning from Expert Solvers}
\label{sec:related-imitation}
% Imitation of an expensive expert (Gasse et al., Paulus et al.); the teacher as a reference row,
% never an upper bound.

% ---------------------------------------------------------------------------------------------
\section{Preliminaries}
\label{sec:prelim}
\label{sec:framework}
% 1.0 page. Definitions only; values and fairness rules go to Section VI.
% Source: EVALUATION_FRAMEWORK_SECTION.md, sections 1-4.

\subsection{Problem Definition}
\label{sec:problem}
% Budgeted, database-wide QDTS. Budget B = 2T + round(s(P - 2T)); report B/P everywhere.

\subsection{Design Requirements}
\label{sec:requirements}
% R1 held-out queries; R2 coverage of query types; R3 comparable budgets; R4 equal information
% access; R5 normalised aggregation; R6 reproducibility and a frozen protocol.

\subsection{Data Split and Budget Accounting}
\label{sec:split}
% Held-out split of trajectories (database vs query pool); overlapping random splits stated
% honestly; user-disjoint splits for learned models.

\subsection{Query Suite}
\label{sec:queries}
% Held-out kNN, Range (point-in-box primary; segment-in-box as a variant), Where-at, held-out join
% over several thresholds; clustering untrained; self-join and self-kNN as diagnostics outside the
% combined score.
% Table: query type, what is compared with what, semantics, score.

\subsection{Scoring}
\label{sec:scoring}
% Endpoints-only floors and normalisation; combined score J4 with per-term reporting;
% recovered share R = (J_method - J_base) / (J_teacher - J_base), R > 1 allowed; statistics
% (mean and std over splits, paired bootstrap intervals, wins per split).

% Training-test separation, information access, external methods, the protocol comparison and the
% threats to validity moved to Section VI (Experimental Setup), which owns all fairness material.

% ---------------------------------------------------------------------------------------------
\section{Structure: When Is QDTS Separable?}
\label{sec:structure}
% 0.75 page. Source: QDTS_THEORY.md. Proofs must fit here in a few lines (no appendix).

\subsection{Stability Under Simplification}
\label{sec:stability}
% Lemma: |d(Q, T_i') - d(Q, T_i)| <= delta_i for original queries; simplification_bound is a
% valid delta_i.

\subsection{Slack: When Query Answers Change}
\label{sec:slack}
% Join, kNN and segment-in-box Range answers are stable while deviation < slack; point-in-box
% Range is not (counterexample).

\subsection{Near-Separability with Original Queries}
\label{sec:separability}
% Each outcome depends on one trajectory; J4 additive up to F1 curvature and kNN near-ties.

\subsection{Optimal Allocation over Concave Envelopes}
\label{sec:allocation-theory}
% Proposition: greedy marginal allocation is exact for concave curves and at envelope vertices
% (Everett; Shoham-Gersho; verify the theorem statements before citing).

\subsection{Where Coupling Remains}
\label{sec:coupling}
% kNN near-ties; compressed (self) queries, where both deviations enter each answer.

% ---------------------------------------------------------------------------------------------
\section{Method: The Teacher and the Student}
\label{sec:method}
% 2.0 pages. The teacher is a teacher or expert, never an oracle or upper bound.

\subsection{Overview}
\label{sec:method-overview}
% Analysis -> teacher -> student; why Section VIII's result makes a per-trajectory student enough.

\subsection{Value Curves and Simplification Recipes}
\label{sec:value-curves}
% Per-trajectory value of query accuracy against budget; the six recipes.

\subsection{Global Budget Allocation}
\label{sec:allocation}
% Greedy allocation along each curve's upper concave envelope; exact budget.

\subsection{Best-Response Loop}
\label{sec:loop}
% Coordinate ascent over recipe choices against exact query trackers; terminates at a local
% optimum.

\subsection{Pairwise Budget Trading}
\label{sec:trading}
% Moves points between trajectories while exact J_T rises. [TR] if pages are short.

\subsection{Complexity}
\label{sec:complexity}
% [TR] if pages are short.

\subsection{The Student}
\label{sec:student}
% Per-trajectory budget regressor and six-way recipe classifier on geometry and workload-density
% features, feeding the same allocator (exact budget guarantee). Design 2 (decision-focused
% value curves) as the route to beating the teacher on held-out join.

% ---------------------------------------------------------------------------------------------
\section{Experimental Setup}
\label{sec:setup}
% About 1.0 page in the conference version; this draft is longer and is meant to be trimmed, with
% the full change lists for RL4QDTS and MLSimp moved to the technical report.
% Scope: how simplifications are evaluated, and why the comparison is fair. Definitions are in
% Section III; runtime, hardware and optimisations are in Section XI.

This section fixes the values behind the definitions of Section~\ref{sec:prelim} and states what
each method may see. Our aim is a comparison in which every method is scored on the same unseen
queries, under the same budget, with the same information available to it wherever the method
allows.

\subsection{Datasets and Splits}
\label{sec:datasets}

We use GeoLife~\todo{cite}. Raw trajectories are cleaned in four steps: repeated timestamps are
removed, trajectories are cut where they leave the Beijing area (more than five consecutive points
outside it), they are split at time gaps longer than 60\,min, and pieces with fewer than 100 points
are dropped. This leaves 19{,}800 trajectories with 19{,}258{,}933 points
(Table~\ref{tab:datasets}). \todo{T-Drive and a third dataset, with their cleaning.}

Each dataset is split at random into a database (85\%) and a query pool (15\%); for GeoLife this
gives 16{,}830 database and 2{,}970 query trajectories. Only the database is simplified, and
the query pool is used for nothing but the test queries. We report five splits (seeds 7, 99, 123,
2024 and 777). A sixth seed (42) is reserved for development and never reported. Because the
splits are drawn from the same trajectories, consecutive databases share most of their content;
the spread across splits therefore measures stability, not independent replication. The student is
trained and tested on splits that are disjoint by GeoLife user, so that repeated routes of one
person cannot appear on both sides \planned{user-disjoint folds}.

\begin{table}[t]
\centering
\caption{Datasets after cleaning.}
\label{tab:datasets}
\small
\begin{tabular}{@{}lrrrr@{}}
\toprule
Dataset & Traj. & Points & Database & Queries \\
\midrule
GeoLife & 19{,}800 & 19{,}258{,}933 & 16{,}830 & 2{,}970 \\
T-Drive & \todo{} & \todo{} & \todo{} & \todo{} \\
\bottomrule
\end{tabular}
\end{table}

\subsection{Budgets}
\label{sec:budgets}

We evaluate four budget shares, $s \in \{0.25, 0.5, 1, 2\}\%$, with the total budget $B$ as defined
in Section~\ref{sec:problem}. Since every trajectory keeps its two endpoints, the share of all
points kept is higher: $B/P = 0.45\%$, $0.70\%$, $1.20\%$ and $2.20\%$, that is, compression by
factors of 220, 142, 83 and 45. All results are reported against $B/P$, which matches the budget
definition of RL4QDTS~\cite{rl4qdts}. For every setting, every method receives the same total $B$,
and endpoints count towards it for every method.

\subsection{Test Workloads}
\label{sec:parameters}

All test workloads are generated from the query pool and the original database of each split, and
are the same for every method (Table~\ref{tab:workloads}).

\begin{itemize}
\item \emph{Held-out kNN:} every query-pool trajectory retrieves its $k = 10$ nearest database
trajectories under continuous Fr\'echet distance; we report recall@10.
\item \emph{Range:} 20{,}000 space-time boxes of $1\,\mathrm{km} \times 1\,\mathrm{km} \times
7$\,days, centred on data points, each with at least 5 answering trajectories (full acceptance from
20), scored by macro F1 with point-in-box semantics. Segment-in-box semantics is reported as a
variant (Section~\ref{sec:shift}).
\item \emph{Where-at:} 128 probe times per database trajectory, restricted to stretches whose
consecutive points are at most 60\,s apart; credit falls linearly from 1 at 25\,m to 0 at 200\,m,
and the share of probes within 50\,m is reported alongside.
\item \emph{Held-out join:} every query-pool trajectory against the database at thresholds of 50,
100, 150 and 200\,m, macro F1 averaged over the four. Three further thresholds (75, 125 and 175\,m),
which no method trains on, test generalisation (Section~\ref{sec:transfer}).
\item \emph{Route clustering:} connected components of size at least 5 at the same four
thresholds. No method optimises it; it is reported but kept out of J4.
\end{itemize}

The self join and self kNN, where both sides are simplified, are reported as a separate score
(J4-self) and never in the headline, because the teacher's training workloads are drawn from the
database itself.

\begin{table}[t]
\centering
\caption{Test workloads per split (GeoLife).}
\label{tab:workloads}
\small
\begin{tabular}{@{}lll@{}}
\toprule
Query & Size & Score \\
\midrule
Held-out kNN & 2{,}970, $k = 10$ & recall@10 \\
Range & 20{,}000 boxes & macro F1 \\
Where-at & 128 per traj. & avg.\ hit; $\le 50$\,m \\
Held-out join & 2{,}970, 4 thresh. & macro F1 \\
Clustering & 4 thresh. & not in J4 \\
\bottomrule
\end{tabular}
\end{table}

\subsection{Metrics and Statistics}
\label{sec:statistics}

The headline score is J4 (Section~\ref{sec:scoring}), always reported next to its four terms, so a
gain on one query type cannot hide a loss on another. The floors are the scores of the
endpoints-only simplification of each split. As a check on the aggregation, we also report the
geometric and harmonic means of the four terms. For each method and budget, we report the mean and
standard deviation over the five splits. Comparisons between two methods are paired, since both are
scored on the same queries; we report the mean paired difference with a 95\% $t$-interval and the
number of splits won. \todo{Bootstrap intervals over queries.} The student is reported by its
recovered share $R$ of the teacher's gain over the best baseline \planned{student results}.

\subsection{Methods Compared}
\label{sec:baselines}

\emph{Error-driven baselines.} RDP, SED and uniform sampling, each with per-trajectory budgets
proportional to $n_i^{\alpha}$ for $\alpha \in \{0, 0.25, 1\}$ (equal, intermediate and
size-proportional), bottom-up PED and DAD \todo{include in the final run}, and database-wide RDP
(\textsc{RDP-Global}), which spends the whole budget by perpendicular error across all trajectories.

\emph{Workload-aware baseline.} \textsc{Partner} gives extra budget, weighted by
$\lambda \in \{0.5, 1\}$, to trajectories that have a join partner in the training workloads.
\todo{Confirm which training pairs it uses.}

\emph{Learned QDTS.} RL4QDTS~\cite{rl4qdts} and MLSimp~\cite{mlsimp}, run as described in
Section~\ref{sec:reproduction}.

\emph{Our methods.} The ablation follows the teacher's stages: allocation with RDP points only
(\textsc{Alloc}), allocation with per-trajectory recipe choice (\textsc{Alloc+Recipe}), plus the
best-response loop (\textsc{+Loop}), plus budget trading (the full teacher), and the student
\planned{student}.

\emph{Tuning rule.} Baselines are reported at every parameter value, and the strongest baseline at
each budget is chosen after seeing the test results. This choice favours the baselines.

% ---------------------------------------------------------------------------------------------
\subsection{Fairness of the Evaluation}
\label{sec:fairness}

A query-driven method sees sample queries and an error-driven one does not, and reproducing
learned methods outside their original code base involves choices. We therefore state the rules
that every method is held to.

\paragraph{Training and test data are separate}
No method sees the query pool. The teacher's training workloads (kNN lists, join pairs, range boxes
and where-at probes) are built from the database alone, with random seeds distinct from those of
the test workloads, and the test workloads are built only after every method has produced its
simplification. All methods are scored on exactly the same queries and probes.

\paragraph{Information access is explicit}
Table~\ref{tab:access} lists what each method sees. Whenever a method uses sample queries, they
come from the training workloads of the same split, never from the test workloads. Where a learned
method can use more training queries than its default, we also report it with the same workload as
the teacher (MLSimp with all 256{,}000 training boxes instead of 100).

\begin{table*}[t]
\centering
\caption{Information available to each method. No method ever sees the held-out query pool.}
\label{tab:access}
\small
\begin{tabular}{@{}p{3.3cm}p{2.6cm}p{6.2cm}p{4.0cm}@{}}
\toprule
Method & Database & Training queries & Tuned on \\
\midrule
Error-driven baselines, \textsc{RDP-Global} & simplified & none & nothing ($\alpha$ reported at every value) \\
\textsc{Partner} & simplified & training join pairs \todo{confirm} & $\lambda$ reported at every value \\
RL4QDTS & simplified & training range boxes of the split \todo{count} & authors' settings, validation on training boxes \\
MLSimp & simplified & 100 training range boxes (also all 256{,}000) & authors' settings and pretrained model as a control \\
Teacher & simplified & kNN lists, join pairs at 4 thresholds, 256{,}000 range boxes, where-at probes, all from the database & development seed and pseudo splits \\
Student & simplified & teacher labels on training users; geometry and density features at inference \planned{} & development seed \\
\bottomrule
\end{tabular}
\end{table*}

\paragraph{Budgets are accounted identically}
Every method receives the same total $B$ with endpoints inside it. MLSimp's code adds endpoints on
top of its budget; we count them inside it, and report the authors' accounting as well.

\paragraph{Learned methods are reproduced faithfully}
\label{sec:reproduction}
Both methods run from the authors' code, included unchanged; every change lives in a separate
adapter and is listed with its reason in our repository \todo{link}. For MLSimp, the 100 range
queries that weight its point scores were its own test queries; we draw them from the training
boxes instead. For RL4QDTS, rewards and validation use the training range boxes of the split, never
the test boxes; a guard checks that every training box finds exactly the trajectories our
evaluator finds. We fix one numerical defect (a hyperbolic tangent that returns NaN for large
inputs), which changes none of the 72{,}678 decisions of the shipped models on our data. For each
method we report the authors' setting, the corrected setting and its strongest variant under our
protocol. MLSimp with the authors' pretrained model scores the same as with ours, so our training
is not the cause of its results. Finally, every method, ours included, is also evaluated under the
authors' native protocol (Table~\ref{tab:protocols}) \todo{native-protocol runs}.

\paragraph{Design choices are fixed before testing}
All design choices of the teacher, including its two query-free terms and their weights, were made
on the development seed or on a pseudo split, which re-splits a database and never reads the query
pool. The teacher weights its six terms equally. The metrics were frozen at commit
\todo{hash}. They had been tightened once, after earlier results were seen; we report the earlier
configuration on the same splits as a check \todo{old-configuration rerun}.

\paragraph{Remaining asymmetries}
Three asymmetries remain, and we report them rather than remove them. First, the teacher is
workload-aware and the error-driven baselines are not; its fair reference is therefore the
workload-aware methods (\textsc{Partner}, RL4QDTS and MLSimp). Second, where-at is held out in time
only: its test probes fall on the same database trajectories the teacher trains on, at other
moments, and range test boxes come from the same generator as the training boxes. These two terms
measure fidelity to the database's own trajectories rather than generalisation to new queries,
which only held-out kNN and join test. Third, point-in-box range rewards where vertices are placed,
not geometric fidelity, which is why we also report segment-in-box semantics.

\begin{table}[t]
\centering
\caption{Evaluation protocols. \todo{Verify every cell for RL4QDTS and MLSimp against the papers.}}
\label{tab:protocols}
\small
\begin{tabular}{@{}llll@{}}
\toprule
 & RL4QDTS & MLSimp & Ours \\
\midrule
DB points & \todo{} & \todo{} & 16.3\,M \\
Test held out & no & no & yes \\
Range boxes & \todo{} & \todo{} & 1\,km, 7\,d \\
Similarity & EDR & \todo{} & Fr\'echet \\
kNN $k$ & 3 & \todo{} & 10 \\
Join thresh. & -- & -- & 50--200\,m \\
Endpoints in $B$ & yes & no & yes \\
\bottomrule
\end{tabular}
\end{table}

% ---------------------------------------------------------------------------------------------
\section{Head-to-Head Comparison}
\label{sec:head-to-head}
% 1.25 pages. Claims C1, C2, C7, C11.

\subsection{Overall Query Accuracy}
\label{sec:overall}
% Figures/tables: F2, T3.

\subsection{Per-Query-Type Results}
\label{sec:per-query}
% Figure: F3. Where the teacher does not win (held-out kNN at 1-2%, held-out join everywhere).

\subsection{Under the Authors' Native Protocol}
\label{sec:native}
% Table: T4.

% ---------------------------------------------------------------------------------------------
\section{What Matters: Decomposing the Gain}
\label{sec:what-matters}
% 1.5 pages. Claims C3, C4, C5.

\subsection{Order-Free Factorial Attribution}
\label{sec:attribution}
% Figures: F4 (and F5 if space).

\subsection{Where the Budget Goes}
\label{sec:budget-distribution}
% Figure: F7.

\subsection{Which Recipes Win}
\label{sec:recipes}
% Figure: F8 (inset in F7 if space).

% ---------------------------------------------------------------------------------------------
\section{When It Matters: Conditions and Limits}
\label{sec:when-matters}
% 1.25 pages. Claims C5, C6, C8, C9. End each subsection with a bold one-line takeaway.

\subsection{The Coordination Gap by Query Type}
\label{sec:coordination-gap}
% Figure: F6.

\subsection{Transfer to Held-Out Queries and the Join Gap}
\label{sec:transfer}
% Figure: F9 (F10 and F15 in [TR]).

\subsection{Range Semantics and Workload Shift}
\label{sec:shift}
% Figure: F11.

\subsection{Robustness of the Ranking}
\label{sec:ranking-robustness}
% One paragraph; table T6 in [TR].

% ---------------------------------------------------------------------------------------------
\section{Student: How Much of the Teacher Can Be Learned?}
\label{sec:student-results}
% 1.0 page. Claims C12, C13 (C14 only if run).

\subsection{Recovered Share and Speed-Up}
\label{sec:recovered-share}
% Table: T10.

\subsection{Generalisation Across Users and Datasets}
\label{sec:student-generalisation}
% Figure: F16.

\subsection{Ablations}
\label{sec:student-ablations}
% Budget head vs budget + recipe; cheap loop pass at the tightest budget; feature groups (T11 in
% [TR]). Cut from the main text first if pages are short.

% ---------------------------------------------------------------------------------------------
\section{Efficiency and Scalability}
\label{sec:efficiency}
% 0.5 page. Claim C10.

\subsection{Simplification Time}
\label{sec:time}
% Figure: F13(a).

\subsection{Scalability with Database Size}
\label{sec:scalability}
% Figure: F13(b).

\subsection{Label Generation versus Inference Cost}
\label{sec:label-cost}
% Teacher label cost against student inference; offline-once argument.

% ---------------------------------------------------------------------------------------------
\section{Limitations and Conclusion}
\label{sec:conclusion}
% 0.25 page.

\subsection{Limitations}
\label{sec:limitations}
% Scope guards; the student inherits the teacher's join overfitting.

\subsection{Conclusion}
\label{sec:conclusion-main}
% Design implications for learned QDTS: per-trajectory scoring priced against one global budget.

% ---------------------------------------------------------------------------------------------
\section*{Acknowledgement of AI-Generated Content}
% Mandatory for ICDE; does not count toward the 12 pages.
\todo{AI-content acknowledgement}

\bibliographystyle{IEEEtran}
% \bibliography{references}

\end{document}
